% A.43: full independent parent proof audit and code rerun passed; registered.
\subsection{An unbounded family of dihedral covers}

\begin{statement}{A.43}{Proved}{Every odd dihedral two-cusp cover has cusp fixed groups at all primes}
For every odd integer \(n\ge3\), let
\[
 C_n:\ v^2=w^{2n+1}-w,
 \qquad X_n=C_n\setminus\{P_0,P_\infty\},
                                                               \tag{A.43.1}
\]
where \(P_0\) and \(P_\infty\) are the unique points
above \(w=0\) and \(w=\infty\). This is a genus-\(n\)
curve and admits a finite etale map \(X_n\to U\) of
degree \(2n\). Geometrically that map is Galois with
dihedral group of order \(2n\). At two its pointed
stable reduction, after finite extension, has a rational
centre carrying both cusps and \(n\) supersingular
elliptic tails; its proper Jacobian has good supersingular
reduction and \([P_0-P_\infty]\) has exact order two.

For every coefficient prime \(\ell\), every cusp packet
at either cusp, and every indicated open derived subgroup,
\[
             \left(\pi_1(X_{n,\bar K})^{(\ell)}\right)^{s(\Gamma_0)}
                                    =I_{P,\ell}.     \tag{A.43.2}
\]
This applies in particular to every induced original packet
and permitted open \(\Delta_0\) of Part A. It proves
Task A.5 for this family of unbounded genus, without a
primality hypothesis on \(n\). It asserts no normalizer
classification for the new higher-genus curves.

The proof includes a general construction. Suppose a smooth
proper positive-genus curve over a number field has good reduction at residue
characteristic \(p\), distinct rational points \(P,Q\)
with the same specialization, and a rational function
\(f\) with \(\operatorname{div}(f)=d(P-Q)\), where
\(d>0\) (exchange \(P,Q\) if necessary).
For any integer \(n\ge2\) with \((n,pd)=1\), the normalized cover \(z^n=f\)
has, after finite extension, a rational component carrying
its two points above \(P,Q\), with \(n\) copies of the
original positive-genus component attached to it. For the
complement of those two points, A.42 gives the same fixed
inertia conclusion at every coefficient prime.
\end{statement}

\begin{proof}
\paragraph{The actual two-branch normalization.}
We first prove the general construction. Work over a finite
local extension at which the proper curve \(C\) has a
smooth model and the distinct marked sections \(P,Q\)
have the same closed point. Separate them by blowing up
until they are disjoint. There is one positive-genus
component \(C_0\), and a chain of rational components
ending in a rational tail carrying both marks. Stabilizing
the pointed model contracts the unmarked valence-two
components and leaves one rational tail \(R_0\) joined
to \(C_0\) by one edge of positive length. The two
marks and that node are three distinct points of \(R_0\).
Use the regular model with the intervening chain retained
for the following divisor calculation.

The horizontal divisor \(dP-dQ\) has degree zero on
every component, because both marked sections meet the
same last rational component. Write \(a_v\) for the
vertical order of \(f\) along component \(v\).
Intersect its principal divisor with each component.
The intersection equations are the graph Laplacian
\[
                       \sum_{u\sim v}(a_u-a_v)=0.
                                                               \tag{A.43.3}
\]
The kernel of the Laplacian on a connected finite graph
consists of constants: multiply each equation by \(a_v\),
sum, and obtain the negative of
\(\sum_{\{u,v\}}(a_u-a_v)^2\). Thus all vertical
orders equal one integer \(a\). If \(\varpi\) is a
uniformizer, \(f_0=\varpi^{-a}f\) has precisely the
horizontal divisor \(dP-dQ\), with no vertical component.
It is a unit at every node. On \(C_0\) its reduction
has divisor zero and is a nonzero constant \(c\); on
all intervening rational components the same is true.
Choose a coordinate \(u\) on the final tail for which
\(P=0\), \(Q=\infty\), and its attaching node is
\(u=1\). There
\[
                         \overline f_0=c u^d.        \tag{A.43.4}
\]
This calculation identifies the reduction of the given
function; it is not a choice of another cover with the
same graph.

Make a finite base extension containing an \(n\)-th
root of \(\varpi^a\), the needed \(n\)-th root of
\(c\), and all \(n\)-th roots of unity. Since
\(n\) is prime to \(p\), the derivative in the
Kummer equation is a unit wherever \(f_0\) is a
unit. On the formal neighbourhood of \(C_0\), the
reduction of \(f_0\) is constant. Hensel's lemma
therefore gives \(n\) separate roots, compatibly
on overlapping formal charts: the root with a fixed
reduction is unique. Thus the actual normalized cover
splits into \(n\) copies of \(C_0\) there. At each
node the complete local ring has equation
\(xy=\varpi^e\), and \(f_0\) is a unit with
constant nonzero reduction. The same Hensel argument
splits the cover into \(n\) copies of this nodal
chart. In particular every connecting annulus is
unramified and keeps its positive length.

On \(R_0\), equation (A.43.4) gives the connected
normalization of
\(z^n=u^d\), after scaling the constant. Because
\((n,d)=1\), choose integers \(a',b'\) with
\(da'+nb'=1\). The parameter
\(s=z^{a'}u^{b'}\) satisfies
\[
                           u=s^n,\qquad z=s^d.
                                                               \tag{A.43.5}
\]
Its smooth compactification is a rational curve. It is
fully ramified at the two marks and unramified at the
node \(u=1\), which has \(n\) distinct preimages.
Near a marked section the same normalized local
parameter gives a smooth chart over the valuation
ring; the ramification is tame. The remaining charts
are etale. Consequently normalization of the actual
model, with the harmless chain components retained,
is semistable. Contracting unmarked valence-two
rational components gives precisely one rational
centre and \(n\) copies of \(C_0\), each attached
at one of those distinct node preimages. The cover
is connected, since valuation \(d\) at a branch
point is coprime to \(n\).

The special dual graph is a tree. Riemann--Hurwitz,
with two totally ramified points, gives
\[
       2g(C')-2=n(2g(C)-2)+2(n-1),
                          \qquad g(C')=n g(C).       \tag{A.43.6}
\]
This agrees with the sum of the genera of its \(n\)
positive-genus components. Its two removed points
are both on the rational centre. Thus A.41--A.42
apply to either selected cusp, proving the general
construction and its all-prime fixed-group assertion.
Additional finite global extensions realizing the
required local extension are harmless, by the field
restriction argument of A.41.

\paragraph{The explicit cover of the once-punctured elliptic curve.}
For the family (A.43.1), put
\[
 E:\ y^2=x^3-x,\qquad T=(0,0),
 \qquad q_n:C_n\longrightarrow E,
          \quad x=w^n,\quad y=w^{(n-1)/2}v.           \tag{A.43.7}
\]
The equation follows on squaring the last expression.
It is the normalization of the cyclic cover \(w^n=x\).
On \(E\),
\(\operatorname{div}(x)=2T-2O\). Since \(n\)
is odd, \(q_n\) has degree \(n\), with total
ramification precisely at \(T,O\). It has no other
ramification. The polynomial defining \(C_n\) has
\(2n+1\) simple roots, and its hyperelliptic model
has genus \(n\), as also follows from (A.43.6).

Over \(K=\mathbf Q(i)\), multiplication by \(1+i\)
on \(E\) is
\[
 [1+i](x,y)=
 \left(\frac{x-x^{-1}}{2i},
             -\frac{1+i}{4}\,y(1+x^{-2})\right).
                                                               \tag{A.43.8}
\]
Its kernel is \(\{O,T\}\). Hence \([1+i]\circ q_n\)
has degree \(2n\) and is etale after precisely the
two points \(P_0,P_\infty\) are deleted and the
point \(O\) is deleted downstairs. This is the
claimed finite etale map \(X_n\to U\).

For the dihedral assertion, after adjoining
\(\zeta_n\) the cyclic generator is
\[
 r(w,v)=(\zeta_n w,\zeta_n^{(n+1)/2}v).
\]
Translation by \(T\) on \(E\), which is
\((x,y)\mapsto(-x^{-1},y x^{-2})\), lifts to
\[
 h(w,v)=\left(-w^{-1},
           (-1)^{(n-1)/2}\frac{v}{w^{n+1}}\right).
                                                               \tag{A.43.9}
\]
Direct substitution gives \(h^2=1\) and
\(hrh=r^{-1}\). These transformations form a
faithful group of order \(2n\) and preserve
\([1+i]\circ q_n\). Since that map already has
that degree, it is the geometric Galois cover with
this dihedral group. In particular this assertion
holds for composite odd \(n\) as well.

\paragraph{Its actual stable reduction at two.}
The elliptic curve \(E\) acquires good supersingular
reduction at two, as in A.40. Its point \(T\) is
nonzero two-torsion, so specializes to \(O\): the
geometric special fibre has no nonzero two-primary
torsion. Apply the general normalization argument
with \(f=x\), \(d=2\), and this odd \(n\).
It gives \(n\) copies of the supersingular elliptic
special fibre attached to one rational centre, with
both marked points on that centre. The Jacobian is
therefore an abelian scheme with special fibre the
product of those \(n\) elliptic curves, by the
compact-type product recalled in
\cite[Section 1]{fghDupuyRabinoff}.
In particular its special \(2\)-rank is zero.
The proper Jacobian has dimension \(n\) and
supersingular reduction.

Finally, on \(C_n\) the function \(w\) has
\[
                  \operatorname{div}(w)=2P_0-2P_\infty.
                                                               \tag{A.43.10}
\]
The divisor class has exact order two: if the two
points were linearly equivalent, a degree-one function
would make the positive-genus curve rational.
Thus both the direct rational-path condition and
the torsion-divisor condition of A.41 are verified
on this actual family. Applying A.42 yields (A.43.2)
for all coefficient primes, including two and every
prime dividing \(n\).

The program \texttt{code/a\_odd\_dihedral\_family.py}
checks the isogeny and cover identities exactly,
the involution and dihedral exponents for sample
odd composite and prime values, the normalized
rational-tail parameter, and the connected graph
Laplacian. The model and all-prime arguments are
proved uniformly above.
\end{proof}
